\documentclass[11pt]{article}

\usepackage[margin=1in]{geometry}
\usepackage{amsmath}
\usepackage{booktabs}
\usepackage{array}

\title{GE-79 Machine Learning Formula Reference}
\author{AI4ALL Diabetes ML Project}
\date{}

\begin{document}

\maketitle

\section*{Purpose}

This document is a human reference for the main classification, model, regression,
and explainable AI formulas discussed in the GE-79 project. The GE-79 project is a
classification project, so regression metrics are included for reference only.

\section*{1. Accuracy}

Accuracy measures the percentage of all predictions that were correct.

\[
\text{Accuracy} =
\frac{TP + TN}{TP + TN + FP + FN}
\]

\noindent Range: \(0\) to \(1\). Higher is better.

\section*{2. Precision}

Precision answers: of everything predicted positive, how many were actually positive?

\[
\text{Precision} =
\frac{TP}{TP + FP}
\]

\noindent Higher precision means fewer false positives.

\section*{3. Recall, or Sensitivity}

Recall answers: of all actual positives, how many were detected?

\[
\text{Recall} =
\frac{TP}{TP + FN}
\]

\noindent Higher recall means fewer false negatives.

\section*{4. F1 Score}

The F1 score balances precision and recall using their harmonic mean.

\[
F_1 =
2 \cdot \frac{\text{Precision} \cdot \text{Recall}}
{\text{Precision} + \text{Recall}}
\]

\noindent Equivalent form using confusion matrix values:

\[
F_1 =
\frac{2TP}{2TP + FP + FN}
\]

\section*{5. Confusion Matrix}

The confusion matrix is not a formula. It is a table that counts correct and
incorrect predictions.

\begin{center}
\begin{tabular}{lcc}
\toprule
 & Predicted Positive & Predicted Negative \\
\midrule
Actual Positive & True Positive (TP) & False Negative (FN) \\
Actual Negative & False Positive (FP) & True Negative (TN) \\
\bottomrule
\end{tabular}
\end{center}

\noindent Every binary classification metric above is computed from these four
values.

\section*{6. Cross-Validation}

For \(k\)-fold cross-validation, the overall score is the average score across
the \(k\) folds:

\[
CV =
\frac{1}{k}\sum_{i=1}^{k} S_i
\]

\noindent For 5-fold cross-validation:

\[
CV =
\frac{S_1 + S_2 + S_3 + S_4 + S_5}{5}
\]

\noindent where \(S_i\) is the metric score, such as accuracy or F1, from fold \(i\).

\section*{7. Random Forest Feature Importance}

Scikit-learn's default Random Forest feature importance is impurity-based. For a
feature \(j\), the unnormalized importance is the total weighted impurity decrease
from all nodes that split on feature \(j\):

\[
FI_j =
\sum_{t \in T_j}
\frac{N_t}{N}
\left(
I_t
- \frac{N_{t_L}}{N_t} I_{t_L}
- \frac{N_{t_R}}{N_t} I_{t_R}
\right)
\]

\noindent where:

\begin{itemize}
    \item \(T_j\) is the set of tree nodes that split on feature \(j\).
    \item \(N\) is the total number of samples.
    \item \(N_t\) is the number of samples at node \(t\).
    \item \(t_L\) and \(t_R\) are the left and right child nodes.
    \item \(I_t\), \(I_{t_L}\), and \(I_{t_R}\) are impurity values.
\end{itemize}

\noindent The normalized feature importances satisfy:

\[
\sum_{j=1}^{p} FI_j = 1
\]

\noindent where \(p\) is the number of features.

\section*{8. Logistic Regression}

Logistic regression predicts the probability of class \(1\) using the sigmoid
function:

\[
P(y = 1 \mid x) =
\frac{1}{1 + e^{-z}}
\]

\noindent where:

\[
z =
\beta_0 + \beta_1 x_1 + \beta_2 x_2 + \cdots + \beta_n x_n
\]

\noindent This is the core equation behind Model 1.

\section*{9. Decision Tree Information Gain}

Decision trees choose splits that improve node purity. With entropy, information
gain is:

\[
IG =
H(\text{parent})
- \sum_{i=1}^{m}
\frac{N_i}{N}
H(\text{child}_i)
\]

\noindent where:

\begin{itemize}
    \item \(H\) is entropy.
    \item \(m\) is the number of child nodes after the split.
    \item \(N_i\) is the number of samples in child node \(i\).
    \item \(N\) is the number of samples in the parent node.
\end{itemize}

\noindent Entropy is:

\[
H =
-\sum_{i=1}^{k} p_i \log_2(p_i)
\]

\noindent where \(k\) is the number of classes and \(p_i\) is the proportion of
samples in class \(i\).

\section*{10. Random Forest Prediction}

For classification, a Random Forest prediction is based on majority vote across
the trees:

\[
\hat{y} =
\operatorname{mode}(T_1(x), T_2(x), \ldots, T_n(x))
\]

\noindent where \(T_i(x)\) is the prediction from tree \(i\) for input \(x\).

\section*{11. Mean Squared Error}

Mean Squared Error is a regression metric and is included for reference only.

\[
MSE =
\frac{1}{n}
\sum_{i=1}^{n}
(y_i - \hat{y}_i)^2
\]

\section*{12. Root Mean Squared Error}

\[
RMSE =
\sqrt{MSE}
=
\sqrt{
\frac{1}{n}
\sum_{i=1}^{n}
(y_i - \hat{y}_i)^2
}
\]

\section*{13. Mean Absolute Error}

\[
MAE =
\frac{1}{n}
\sum_{i=1}^{n}
\left|y_i - \hat{y}_i\right|
\]

\section*{14. Coefficient of Determination, \(R^2\)}

\[
R^2 =
1 - \frac{SS_{res}}{SS_{tot}}
\]

\noindent where:

\[
SS_{res} =
\sum_{i=1}^{n}
(y_i - \hat{y}_i)^2
\]

\[
SS_{tot} =
\sum_{i=1}^{n}
(y_i - \bar{y})^2
\]

\section*{15. Explainable AI: SHAP}

SHAP expresses a model prediction as a baseline plus feature-level contributions:

\[
f(x) =
\phi_0 +
\sum_{i=1}^{M}
\phi_i
\]

\noindent where:

\begin{itemize}
    \item \(\phi_0\) is the baseline prediction.
    \item \(\phi_i\) is the SHAP contribution from feature \(i\).
    \item \(M\) is the number of features.
\end{itemize}

\noindent Each SHAP value explains how much a feature pushes the prediction above
or below the baseline.

\section*{16. Gini Impurity}

Gini impurity is a common splitting criterion for decision trees and random
forests:

\[
G =
1 - \sum_{i=1}^{k} p_i^2
\]

\noindent where:

\begin{itemize}
    \item \(k\) is the number of classes.
    \item \(p_i\) is the probability, or node proportion, of class \(i\).
\end{itemize}

\noindent Lower Gini impurity indicates a purer node.

\end{document}
